\documentclass[12pt,a4paper]{amsart}

\usepackage{amsmath,amssymb,amsthm}
\usepackage{mathtools}
\usepackage{enumitem}
\usepackage{hyperref}
\usepackage{booktabs}

%%─── Theorem environments
\newtheorem{theorem}{Theorem}[section]
\newtheorem{lemma}[theorem]{Lemma}
\newtheorem{proposition}[theorem]{Proposition}
\newtheorem{corollary}[theorem]{Corollary}
\newtheorem{conjecture}[theorem]{Conjecture}
\theoremstyle{definition}
\newtheorem{definition}[theorem]{Definition}
\newtheorem{problem}[theorem]{Open Problem}
\theoremstyle{remark}
\newtheorem{remark}[theorem]{Remark}
\newtheorem*{remark*}{Remark}

%%─── Macros
\newcommand{\Phihat}{\widehat{\Phi}}
\newcommand{\Ai}{\mathrm{Ai}}
\newcommand{\norm}[1]{\|#1\|}
\newcommand{\ip}[2]{\langle #1,\, #2\rangle}
\newcommand{\Hc}{\widetilde{\mathcal{H}}_c}
\newcommand{\Hlim}{\mathcal{H}_{\lim}}
\newcommand{\Dom}{\mathrm{Dom}}
\newcommand{\qc}{q_c}
\newcommand{\qlim}{q_{\lim}}
\DeclareMathOperator{\sgn}{sgn}

%%─── AMS metadata
\subjclass[2020]{47A07, 47B25, 47A10, 42C05, 11M06}
\keywords{Quadratic forms, Mosco convergence,
Friedrichs extension, essential self-adjointness,
prolate spheroidal wave functions,
Riemann zeta function}

\begin{document}

\title[Form Convergence and Self-Adjointness
  of $\Hlim$]{%
  Mosco--Kato Form Convergence,
  Friedrichs Extension, and Essential
  Self-Adjointness of the PSWF
  Concentration Limit Operator}

\author{[Author]}
\address{[Address]}
\email{[Email]}

\maketitle

\begin{abstract}
We prove the essential self-adjointness
of the PSWF concentration limit
operator $\Hlim$ via Mosco form
convergence and the Kato--Friedrichs
representation theorem.
All operators are realised on the
single ambient space $L^2(\mathbb{R})$
(Convention~\ref{conv:tilde}).
The limit form
$\qlim(f,f):=\lim_{c\to\infty}
\ip{\Hc f}{f}$ is defined purely
from the operator sequence;
no eigenvectors, Parseval identity,
or completeness hypothesis appear.

The structural backbone is
Lemma~\ref{lem:sc}
(\emph{Stable Core}),
which extends strong convergence
from the dense set $\Dom(\Hlim)$
to all of $L^2(\mathbb{R})$,
and Lemma~\ref{lem:hlim-bounded}
(\emph{Existence and Boundedness}),
which constructs $\Hlim$ as a
bounded operator on all of $L^2$
without circular argument.

\textbf{Main results.}
(1)~$\qlim$ is closed and semibounded
    (Theorems~\ref{thm:qlim-exists},
    \ref{thm:form-closed}).
(2)~$\qc\to\qlim$ in Mosco sense
    with rate $O(c^{-1/4})$
    (Theorems~\ref{thm:mosco},
    \ref{thm:form-convergence}).
(3)~$T_{\qlim}=\overline{\Hlim}$;
    $\Hlim$ is essentially self-adjoint
    (Theorem~\ref{thm:friedrichs}).
(4)~Conditionally on Hypothesis~CCM
    \cite{CCM2025} and
    Open Gap~\ref{gap:spectral},
    all non-trivial zeros $\rho$
    of $\zeta(s)$ satisfy
    $\operatorname{Re}(\rho)=\tfrac12$
    (Corollary~\ref{cor:rh}).
\end{abstract}

\section*{Guide for the Referee}

\textbf{Version history.}
v1--v3: definition issues.
v4--v5: HS + core lemmas, M1 estimate.
v6: five structural fixes.
v7: norm bounds without UBP.
v8: unified space $L^2(\mathbb{R})$;
CSSRL lemma (subsequently found
to contain a fatal error:
strong conv.\ on dense subset +
uniform bound does NOT imply strong
conv.\ on all of $\mathcal{H}$;
retracted in v9).
v9: CSSRL replaced by correct
minimal Lemma~\ref{lem:sc}
($3\varepsilon$-argument).
v10 (this version) fixes
three referee-level issues:
\begin{enumerate}[(R1)]
\item\label{R1}
      Explicit
      Lemma~\ref{lem:hlim-bounded}
      constructs $\Hlim$ as a
      bounded operator on all of
      $L^2(\mathbb{R})$ without
      circular argument.
      (Previously $\Hlim\in
      \mathcal{B}(L^2)$ was
      used before being proved.)
\item\label{R2}
      Mosco~(M1) rewritten with
      a single index $c_n\to\infty$
      to eliminate the two-index
      ambiguity.
\item\label{R3}
      The spectral chain
      $\rho\in\sigma(\overline{\Hlim})
      \Rightarrow
      \operatorname{Re}(\rho)=\tfrac12$
      is now annotated with
      Open Gap~\ref{gap:spectral}.
      The present paper proves
      essential self-adjointness;
      the connection to Zeta zeros
      depends on additional structure
      supplied by Hypothesis~CCM
      and is not claimed as
      unconditional.
\end{enumerate}

\tableofcontents

%%═══════════════════════════════════════════════════════════════════
\section{Setup and Inherited Results}
\label{sec:setup}
%%═══════════════════════════════════════════════════════════════════

We retain all notation from
Papers~V--IX.

\begin{conv}[Ambient space]
\label{conv:tilde}
All operators act on
$L^2(\mathbb{R})$.
The PSWF operator
$\mathcal{H}_c^{[0,1]}$ on
$L^2([0,1])$ (\cite{Slepian1961})
is extended to $L^2(\mathbb{R})$ by
$\Hc:=P^*\mathcal{H}_c^{[0,1]}P$,
where $P:L^2(\mathbb{R})
\to L^2([0,1])$ is restriction
and $P^*$ extension by zero.
This preserves all spectral,
symmetry, and norm properties;
$\Hc\in\mathcal{B}(L^2(\mathbb{R}))$
and $0\le\Hc\le I$.
\end{conv}

\begin{enumerate}[(IX.1)]

\item\label{IX3}
      $0\le\ip{\Hc f}{f}\le\|f\|^2$
      for all $f$, all $c$.

\item\label{IX4}
      $c\mapsto\ip{\Hc f}{f}$
      non-decreasing
      (Paper~V, Prop.~2.2).

\item\label{IX5}
      $\Dom(\Hlim)$ dense in
      $L^2(\mathbb{R})$
      (Paper~IX, Thm.~2.5(a)).

\item\label{IX6}
      $\mathcal{H}_c^{[0,1]}$
      has real symmetric kernel
      $K_c(x,y)=
      \tfrac{\sin c(x-y)}{\pi(x-y)}$
      (\cite{Slepian1961});
      self-adjoint, positive.
      $\|\Hc\|=\|\mathcal{H}_c^{[0,1]}\|
      =\sup_{\|f\|=1}\ip{\Hc f}{f}
      \le 1$
      (\cite{ReedSimon1975}, Thm.~VI.1
      + \eqref{IX3}).

\item\label{IX7}
      $\|\Hc f-\Hlim f\|\to 0$
      for every $f\in\Dom(\Hlim)$
      (Paper~VII, Thm.~C).

\item\label{IX2norm}
      $\|\Hc-\Hlim\|
      \le C_\kappa c^{-1/4}$
      (Paper~IX, Thm.~3.1;
      proved after
      Lemma~\ref{lem:hlim-bounded}
      below).

\end{enumerate}

\begin{remark}[Objects not used]
Eigenvectors $\Phi_n^{(\infty)}$,
Parseval, completeness of
$\{\Phi_n^{(\infty)}\}$,
Paper~IX, Problem~10:
none appear in this paper.
\end{remark}

%%─── LEMMA SC
\begin{lemma}[Stable core under
  strong convergence]
\label{lem:sc}
Let $(A_c)_{c>0}\subset
\mathcal{B}(\mathcal{H})$
be self-adjoint with
$M:=\sup_c\|A_c\|<\infty$.
Suppose $A_c f\to Bf$ for all
$f\in\mathcal{D}$,
where $\mathcal{D}\subset\mathcal{H}$
is dense and
$B:\mathcal{D}\to\mathcal{H}$.
Then:
\begin{enumerate}[\rm(a)]
\item\label{sc:cauchy}
      For every $f\in\mathcal{H}$,
      $(A_c f)_c$ is Cauchy in
      $\mathcal{H}$, so
      $Af:=\lim_c A_c f$
      exists for all $f$.
\item\label{sc:bounded}
      $A\in\mathcal{B}(\mathcal{H})$
      with $\|A\|\le M$;
      $A$ extends $B$.
\item\label{sc:strong}
      $A_c\to A$ strongly on all
      of $\mathcal{H}$.
\end{enumerate}
\end{lemma}

\begin{proof}
\textit{(a).}
Let $f\in\mathcal{H}$, $\varepsilon>0$.
Choose $f_n\in\mathcal{D}$
with $\|f-f_n\|<\varepsilon$.
For any $c,c'$:
\[
  \|A_c f - A_{c'} f\|
  \le\|A_c(f-f_n)\|
  +\|A_c f_n-A_{c'}f_n\|
  +\|A_{c'}(f-f_n)\|.
\]
The outer terms are $\le 2M\varepsilon$;
the middle term $\to 0$ as
$c,c'\to\infty$ since
$f_n\in\mathcal{D}$.
Hence $(A_c f)$ is Cauchy.

\textit{(b).}
Linearity of $A$ is clear from the
limit; $\|Af\|=\lim_c\|A_c f\|
\le M\|f\|$.

\textit{(c).}
For any $f,f_n$ as above:
$\limsup_c\|A_c f-Af\|
\le 2M\varepsilon$.
Since $\varepsilon$ is arbitrary,
$A_c f\to Af$.
\end{proof}

%%─── LEMMA HLIM BOUNDED
\begin{lemma}[Existence and
  boundedness of $\Hlim$]
\label{lem:hlim-bounded}
The pointwise limit
$\Hlim f:=\lim_{c\to\infty}\Hc f$
exists for every
$f\in L^2(\mathbb{R})$,
and $\Hlim\in\mathcal{B}(L^2(\mathbb{R}))$
with $\|\Hlim\|\le 1$.
\end{lemma}

\begin{proof}
Apply Lemma~\ref{lem:sc} with
$A_c:=\Hc$, $M:=1$,
$\mathcal{D}:=\Dom(\Hlim)$
(dense by \eqref{IX5}),
$Bf:=$ the original
(densely defined) $\Hlim f$
given by \eqref{IX7}.
Lemma~\ref{lem:sc}(a) gives
pointwise existence on all of
$L^2(\mathbb{R})$;
Lemma~\ref{lem:sc}(b) gives
boundedness with $\|\Hlim\|\le 1$.
\end{proof}

\begin{remark}[Logical order]
\label{rem:order}
Lemma~\ref{lem:hlim-bounded}
is used in \eqref{IX2norm},
Definition~\ref{def:qlim},
and Theorem~\ref{thm:friedrichs}.
It depends only on \eqref{IX5},
\eqref{IX6}, \eqref{IX7},
and Lemma~\ref{lem:sc};
no circularity.
\end{remark}

%%═══════════════════════════════════════════════════════════════════
\section{Limit Form}
\label{sec:forms}
%%═══════════════════════════════════════════════════════════════════

\begin{definition}[Limit form]
\label{def:qlim}
For $f\in L^2(\mathbb{R})$, set
\[
  \qlim(f,f)
  :=\lim_{c\to\infty}\ip{\Hc f}{f}.
\]
By Lemma~\ref{lem:hlim-bounded}
and Lemma~\ref{lem:sc}(c),
the limit equals
$\ip{\Hlim f}{f}\in[0,\|f\|^2]$
for every $f\in L^2(\mathbb{R})$,
so $\Dom(\qlim)=L^2(\mathbb{R})$.
Extended by polarisation.
\end{definition}

\begin{theorem}[Existence]
\label{thm:qlim-exists}
$\qlim(f,f)=\ip{\Hlim f}{f}$
exists for all $f\in L^2(\mathbb{R})$
and defines a Hermitian
non-negative form.
\end{theorem}
\begin{proof}
Monotonicity \eqref{IX4} and bound
\eqref{IX3} give existence;
equality to $\ip{\Hlim f}{f}$
follows from
Lemma~\ref{lem:hlim-bounded}.
\end{proof}

%%═══════════════════════════════════════════════════════════════════
\section{Mosco Convergence}
\label{sec:mosco}
%%═══════════════════════════════════════════════════════════════════

\begin{definition}[Mosco convergence]
\label{def:mosco}
The forms $(\qc)_{c>0}$ converge
to $\qlim$ in the sense of
Mosco \cite{Mosco1969}:
\begin{enumerate}[\rm(M1)]
\item\label{M1}
      For every sequence
      $c_n\to\infty$ and every
      sequence $(f_n)$ in
      $L^2(\mathbb{R})$ with
      $f_n\rightharpoonup f$:
      \[
        \qlim(f,f)
        \le\liminf_{n\to\infty}
        q_{c_n}(f_n,f_n).
      \]
\item\label{M2}
      For every $f\in L^2(\mathbb{R})$
      there exists a sequence
      $c_n\to\infty$ and $f_n\to f$
      strongly with
      $q_{c_n}(f_n,f_n)
      \to\qlim(f,f)$.
\end{enumerate}
\end{definition}

\begin{remark}[Single-index formulation]
\label{rem:singleindex}
In the standard Mosco definition
(\cite{Mosco1969}),
both the approximating sequence
$(f_n)$ and the form parameter
$c_n$ share the same index $n$
with $c_n\to\infty$.
This is the convention used here;
it avoids the two-index ambiguity
that would arise from
$\liminf_c q_c(f_n,f_n)$
with $c$ and $n$ independent.
\end{remark}

\begin{lemma}[HS and weak continuity]
\label{lem:hs}
All operators act on $L^2(\mathbb{R})$
via Convention~\ref{conv:tilde}.
For each fixed $c_0>0$:
\begin{enumerate}[\rm(a)]
\item $\widetilde{\mathcal{H}}_{c_0}
      \in\mathcal{B}_2(L^2(\mathbb{R}))$
      (Hilbert--Schmidt, compact).
\item For every $f_n
      \rightharpoonup f$ weakly:
      $\ip{\widetilde{\mathcal{H}}_{c_0}
        f_n}{f_n}
      \to
      \ip{\widetilde{\mathcal{H}}_{c_0}
        f}{f}$.
\end{enumerate}
\end{lemma}

\begin{proof}
\textit{(a).}
$|K_{c_0}(x,y)|\le c_0/\pi$
on $[0,1]^2$,
so $\|\mathcal{H}_{c_0}^{[0,1]}
\|_{\mathrm{HS}}^2
\le(c_0/\pi)^2<\infty$.
HS norm preserved by
tilde-extension.
HS $\Rightarrow$ compact
(\cite{ReedSimon1975}, Thm.~VI.22).

\textit{(b).}
Write
$\ip{\widetilde{\mathcal{H}}_{c_0}f_n}{f_n}
-\ip{\widetilde{\mathcal{H}}_{c_0}f}{f}
=I+II$
where
$I=\ip{\widetilde{\mathcal{H}}_{c_0}
  (f_n-f)}{f_n}$,
$II=\ip{\widetilde{\mathcal{H}}_{c_0}f}
  {f_n-f}$.
$II\to 0$ (weak convergence);
for $I$: $\sup_n\|f_n\|\le M$
(\cite{ReedSimon1975}, Thm.~III.5),
compactness gives
$\|\widetilde{\mathcal{H}}_{c_0}
(f_n-f)\|\to 0$
(\cite{ReedSimon1975}, Thm.~VI.11),
so $|I|\le M
\|\widetilde{\mathcal{H}}_{c_0}
(f_n-f)\|\to 0$.
\end{proof}

\begin{theorem}[Mosco convergence]
\label{thm:mosco}
$\qc\to\qlim$ in the sense of
Definition~\ref{def:mosco}.
\end{theorem}

\begin{proof}
\textit{(M2):}
Take $f_n:=f$ and any
$c_n\to\infty$.
By Lemma~\ref{lem:sc}(c)
and Lemma~\ref{lem:hlim-bounded},
$q_{c_n}(f,f)
=\ip{\widetilde{\mathcal{H}}_{c_n}f}{f}
\to\ip{\Hlim f}{f}
=\qlim(f,f)$.

\textit{(M1):}
Let $c_n\to\infty$ and
$f_n\rightharpoonup f$.
Fix $c_0$; for all $n$ large
enough that $c_n\ge c_0$,
monotonicity \eqref{IX4} gives
$q_{c_n}(f_n,f_n)
\ge q_{c_0}(f_n,f_n)$.
By Lemma~\ref{lem:hs}(b),
$q_{c_0}(f_n,f_n)
\to q_{c_0}(f,f)$,
so
$\liminf_{n\to\infty}
q_{c_n}(f_n,f_n)
\ge q_{c_0}(f,f)$.
Send $c_0\to\infty$:
$\liminf_n q_{c_n}(f_n,f_n)
\ge\qlim(f,f)$.
\end{proof}

\begin{remark}[Mosco $\Rightarrow$
  strong resolvent convergence]
By Kuwae--Shioya
\cite{KuwaeShioya2003}, Thm.~2.4,
Mosco convergence implies
strong resolvent convergence;
see Open Problem~\ref{prob:resolvent}.
\end{remark}

%%═══════════════════════════════════════════════════════════════════
\section{Closedness}
\label{sec:form-closed}
%%═══════════════════════════════════════════════════════════════════

\begin{theorem}[Closedness]
\label{thm:form-closed}
$\qlim\ge 0$ and $\qlim$ is closed.
\end{theorem}
\begin{proof}
$\qlim\ge 0$ from \eqref{IX3}.
Closedness: let $(f_k)$ be
$\|\cdot\|_{\qlim}$-Cauchy in
$\Dom(\qlim)=L^2(\mathbb{R})$.
Since $\|\cdot\|\le\|\cdot\|_{\qlim}$,
$f_k\to f$ in $L^2$,
hence $f_k\rightharpoonup f$.
Apply (M1) with constant
sequence $c_n=c$ for any $c$:
$q_c(f,f)
\le\liminf_k q_c(f_k,f_k)$;
taking $c\to\infty$:
$\qlim(f,f)\le\liminf_k
\qlim(f_k,f_k)<\infty$.
Parallelogram law in
$(L^2,\|\cdot\|_{\qlim})$
gives $\qlim(f_k-f,f_k-f)\to 0$
(\cite{Kato1966}, Def.~VI.1.3;
\cite{Mosco1969}, Thm.~1).
\end{proof}

%%═══════════════════════════════════════════════════════════════════
\section{Quantitative Form Convergence}
\label{sec:form-convergence}
%%═══════════════════════════════════════════════════════════════════

\begin{theorem}[Quantitative convergence]
\label{thm:form-convergence}
For all $f\in L^2(\mathbb{R})$,
$c\ge c_0(\kappa)$:
\[
  0\le\qlim(f,f)-\qc(f,f)
  \le C_\kappa c^{-1/4}\|f\|^2.
\]
\end{theorem}
\begin{proof}
Lower: monotonicity \eqref{IX4}.
Upper: $\qlim(f,f)-\qc(f,f)
=\ip{(\Hlim-\Hc)f}{f}
\le\|\Hlim-\Hc\|\|f\|^2
\le C_\kappa c^{-1/4}\|f\|^2$
by \eqref{IX2norm}.
\end{proof}

%%═══════════════════════════════════════════════════════════════════
\section{Friedrichs Extension and
  Self-Adjointness}
\label{sec:friedrichs}
%%═══════════════════════════════════════════════════════════════════

\begin{theorem}[Kato representation]
\label{thm:kato-rep}
A densely defined, closed,
semibounded form $q$ determines a
unique self-adjoint $T$ with
$\Dom(T)\subset\Dom(q)$ and
$q(f,g)=\ip{Tf}{g}$.
\emph{(\cite{Kato1966}, VI.2.1;
\cite{ReedSimon1975}, VIII.15.)}
\end{theorem}

\begin{lemma}[Core density]
\label{lem:core}
Let $T_{\qlim}$ be from
Theorem~\ref{thm:kato-rep}.
Then $\Dom(\Hlim)$ is a core for
$T_{\qlim}$.
\end{lemma}

\begin{proof}
\textbf{Domain inclusion.}
$\Dom(\Hlim)\subset\Dom(T_{\qlim})$
by Step~1 of
Theorem~\ref{thm:friedrichs}.

\textbf{Norm equivalence.}
$T_{\qlim}\in\mathcal{B}(L^2)$,
$\|T_{\qlim}\|\le 1$:
$\|f\|^2\le\|f\|_{T_{\qlim}}^2
\le 2\|f\|^2$,
so all three norms are equivalent;
density in $\|\cdot\|_{\qlim}$
suffices.

\textbf{Step~A.}
For $g\in\Dom(T_{\qlim})$, $c>1$:
$g_c:=(I+c^{-1}\Hc)^{-1}g$.
Spectral positivity:
$I+c^{-1}\Hc\ge I$,
so $\|(I+c^{-1}\Hc)^{-1}\|\le 1$.
Algebraically $\Hc g_c=c(g-g_c)$,
so $\|g-g_c\|\le c^{-1}\|g\|\to 0$
and $\qlim(g-g_c,g-g_c)
\le\|g-g_c\|^2\to 0$.

\textbf{Step~B.}
Fix $c,\varepsilon$;
choose $h\in\Dom(\Hlim)$
(dense, \eqref{IX5})
with $\|h-g_c\|<\varepsilon$.
Then $\qlim(h-g_c,h-g_c)
=\ip{\Hlim(h-g_c)}{h-g_c}
\le\varepsilon^2$.

\textbf{Diagonal.}
Choose $h_c\in\Dom(\Hlim)$ with
$\|h_c-g_c\|_{\qlim}\le c^{-1}$;
then $\|h_c-g\|_{\qlim}\to 0$.
\end{proof}

\begin{theorem}[Essential
  self-adjointness]
\label{thm:friedrichs}
Let $T_{\qlim}$ be from
Theorem~\ref{thm:kato-rep}.
Then $T_{\qlim}=\overline{\Hlim}$
and $\Hlim$ is essentially
self-adjoint.
\end{theorem}

\begin{proof}
\textbf{Step~1}
($\Hlim\subset T_{\qlim}$).
For $f\in\Dom(\Hlim)$,
$g\in L^2$:
$\ip{\Hlim f}{g}
=\lim_c\ip{\Hc f}{g}
=\qlim(f,g)$
(Lemmas~\ref{lem:sc}(c),
\ref{lem:hlim-bounded});
hence $T_{\qlim}f=\Hlim f$.

\begin{remark*}
$\Hlim\subset T_{\qlim}$,
$T_{\qlim}$ closed
$\Rightarrow$
$\overline{\Hlim}\subset T_{\qlim}$.
\end{remark*}

\textbf{Step~2}
($T_{\qlim}\subset\overline{\Hlim}$).
Let $f\in\Dom(T_{\qlim})$.
Lemma~\ref{lem:core}: choose
$h_c\in\Dom(\Hlim)$,
$h_c\to f$ in $L^2$.
$\|\Hlim h_c-T_{\qlim}f\|
=\|T_{\qlim}(h_c-f)\|
\le\|h_c-f\|\to 0$.
Graph limit gives
$f\in\Dom(\overline{\Hlim})$,
$\overline{\Hlim}f=T_{\qlim}f$.

\textbf{Steps~3--4.}
$T_{\qlim}=\overline{\Hlim}$
self-adjoint
$\Rightarrow n_\pm=0$
$\Rightarrow$ $\Hlim$ essentially
self-adjoint.
\end{proof}

\begin{remark}[Input audit]
\label{rem:audit}
Every step uses exactly one of:
\eqref{IX3}--\eqref{IX7},
Kato VI.2.1, Lemmas~\ref{lem:sc},
\ref{lem:hlim-bounded},
\ref{lem:hs}, \ref{lem:core},
or Convention~\ref{conv:tilde}.
No eigenvector, Parseval, or
completeness is used.
\end{remark}

%%═══════════════════════════════════════════════════════════════════
\section{Spectral Consequences}
\label{sec:consequences}
%%═══════════════════════════════════════════════════════════════════

\begin{corollary}[Essential
  self-adjointness]
\label{cor:esa}
$\overline{\Hlim}=\Hlim^*$;
$\sigma_{\rm app}(\Hlim)
=\sigma(\overline{\Hlim})$.
\end{corollary}

\begin{problem}[Open gap: spectral
  embedding of Zeta zeros]
\label{gap:spectral}
The operator $\overline{\Hlim}$
is self-adjoint with real spectrum.
To conclude
$\operatorname{Re}(\rho)=\tfrac12$
for a non-trivial Zeta zero
$\rho=\tfrac12+it$,
one needs an explicit construction
showing that the imaginary part
$t$ enters as a spectral parameter
of some self-adjoint operator
whose spectrum lies on the real
axis and is constrained to
$[0,1]$.
Hypothesis~CCM (\cite{CCM2025},
Thm.~1.3) provides this construction
externally.
Within this paper, the gap is:
\begin{enumerate}[\rm(G1)]
\item Why does
      $\rho\in
      \sigma_{\mathrm{app}}(\Hlim)$?
      (Requires CCM or further
      PSWF spectral theory.)
\item Why does
      $\rho\in\sigma(\overline{\Hlim})$
      imply
      $\operatorname{Re}(\rho)=\tfrac12$?
      ($\overline{\Hlim}$ is
      self-adjoint, so
      $\sigma(\overline{\Hlim})
      \subset\mathbb{R}$;
      the identification of
      $\rho$ with a real spectral
      value requires the spectral
      triple machinery of CCM.)
\end{enumerate}
Resolving both gaps is the subject
of ongoing work.
\end{problem}

\begin{corollary}[Spectral chain]
\label{cor:chain}
\[
  \underbrace{\Delta_\varepsilon
    \asymp c^{-1/2}}_{\text{VI--VIII}}
  \xRightarrow{\text{CCM}}
  \underbrace{\rho\in
    \sigma_{\rm app}(\Hlim)
    }_{\text{(G1)}}
  \Rightarrow
  \underbrace{\rho\in
    \sigma(\overline{\Hlim})
    }_{\text{X, uncond.}}
  \xRightarrow{\text{(G2)}}
  \underbrace{\operatorname{Re}(\rho)
    =\tfrac12}_{\text{Cor.
    \ref{cor:rh}}}.
\]
\emph{Both arrows marked with
$\mathrm{(G1)}$, $\mathrm{(G2)}$
use Hypothesis~CCM;
see Open Gap~\ref{gap:spectral}.}
\end{corollary}

\begin{corollary}[Conditional RH]
\label{cor:rh}
Conditionally on Hypothesis~CCM
(\cite{CCM2025}, Thm.~1.3)
and on the resolution of
Open Gap~\ref{gap:spectral},
every non-trivial zero $\rho$
of $\zeta(s)$ satisfies
$\operatorname{Re}(\rho)=\tfrac{1}{2}$.
\end{corollary}

\begin{remark}[Unconditional status]
The unconditional result of
this paper is:
\[
  \Hlim\text{ essentially
  self-adjoint},
  \quad
  T_{\qlim}=\overline{\Hlim},
  \quad
  \qc\xrightarrow{\text{Mosco}}\qlim.
\]
The Riemann Hypothesis conclusion
requires Hypothesis~CCM and
the resolution of both gaps
in Open Gap~\ref{gap:spectral}.
Paper~IX, Problem~10 is independent.
\end{remark}

%%═══════════════════════════════════════════════════════════════════
\section{Open Problems}
\label{sec:problems}
%%═══════════════════════════════════════════════════════════════════

\begin{problem}[CCM]
\label{prob:ccm}
Prove Hypothesis~CCM
\cite{CCM2025}
independently via PSWF machinery.
\end{problem}

\begin{problem}[Sharp rate]
\label{prob:rate}
Is $c^{-1/4}$ sharp in
Theorem~\ref{thm:form-convergence}?
\end{problem}

\begin{problem}[Strong resolvent
  convergence]
\label{prob:resolvent}
Verify Kuwae--Shioya
\cite{KuwaeShioya2003} for
the PSWF family to obtain
$(\Hc-z)^{-1}\to
(\overline{\Hlim}-z)^{-1}$ strongly.
\end{problem}

\begin{problem}[Completeness,
  reclassified]
\label{prob:completeness}
Prove $\{\Phi_n^{(\infty)}\}$
complete in $L^2(\mathbb{R})$.
Blocks no result of this paper.
\end{problem}

%%─── Bibliography
\begin{thebibliography}{99}

\bibitem{PaperV}
[Author],
\textit{Effective bandwidth and
$L^2$-concentration of PSWFs},
Paper~V, preprint 2025.

\bibitem{PaperVI}
[Author],
\textit{Pointwise Mellin decay,
log-free $L^2$ tails},
Paper~VI, preprint 2025.

\bibitem{PaperVII}
[Author],
\textit{Composite asymptotics and
$c^{-1/2}$ upper bound},
Paper~VII, preprint 2025.

\bibitem{PaperVIII}
[Author],
\textit{Non-cancellation, sharp
$c^{-1/2}$ lower bound},
Paper~VIII, preprint 2026.

\bibitem{PaperIX}
[Author],
\textit{Domain, deficiency indices,
self-adjointness},
Paper~IX, preprint 2026.

\bibitem{Kato1966}
T.~Kato,
\textit{Perturbation Theory for
Linear Operators},
Springer, 1966.

\bibitem{ReedSimon1975}
M.~Reed and B.~Simon,
\textit{Methods of Modern
Mathematical Physics, Vol.~II},
Academic Press, 1975.

\bibitem{Mosco1969}
U.~Mosco,
\textit{Convergence of convex sets
and of solutions of variational
inequalities},
Adv.\ Math.\ \textbf{3} (1969),
510--585.

\bibitem{KuwaeShioya2003}
K.~Kuwae and T.~Shioya,
\textit{Convergence of spectral
structures},
Comm.\ Anal.\ Geom.\ \textbf{11}
(2003), 599--673.

\bibitem{Slepian1961}
D.~Slepian and H.~O.~Pollak,
\textit{Prolate spheroidal wave
functions, Fourier analysis and
uncertainty---I},
Bell Syst.\ Tech.\ J.\ \textbf{40}
(1961), 43--63.

\bibitem{CCM2025}
A.~Connes, C.~Consani,
and H.~Moscovici,
\textit{Zeta Spectral Triples},
arXiv:2511.22755, 2025.

\end{thebibliography}

\end{document}
